\documentclass[10pt,a4paper]{amsart}
\usepackage[margin=19mm]{geometry}
\usepackage{amsmath,amssymb,mathtools}
\usepackage[scr=rsfs,cal=euler]{mathalfa}
\usepackage{xfrac}
\usepackage[unicode,hidelinks,bookmarks=false]{hyperref}
\newcommand{\ie}{i.e.\ }
\newcommand{\cf}{cf.\ }
\newcommand{\resp}{respectively\ }
\newcommand{\Subsection}{\subsection}
\providecommand{\nobreakdash}{\nobreak\hbox{-}}
\setlength{\emergencystretch}{2em}
\multlinegap0pt
\usepackage{bm}
\usepackage{bbm}
\usepackage{dsfont}
\usepackage{xspace}
\usepackage{cancel}
\usepackage{mathtools}
\usepackage{cleveref}
\usepackage{amsthm}
\makeatletter
\@ifundefined{theorem}{\newtheorem{theorem}{Theorem}}{}
\makeatother
\title[Comparison Theorems for the Mixing Times of Systematic and R]{Comparison Theorems for the Mixing Times of Systematic and Random Scan Dynamics}
\author{Jason Gaitonde, Elchanan Mossel}
\date{Source: arXiv:2410.11136v2 (CC BY 4.0)}
\begin{document}
\maketitle

\noindent\emph{Source and licence.} Comparison Theorems for the Mixing Times of Systematic and Random Scan Dynamics. Version arXiv:2410.11136v2 (CC BY 4.0), distributed under the Creative Commons Attribution 4.0 International licence, as recorded in the arXiv licence metadata for this exact version and shown on the abstract page https://arxiv.org/abs/2410.11136v2. Full rights notice: licenses/arxiv-2410.11136-CC-BY-4.0.txt.\par\smallskip

\noindent\textit{Editorial scope.} Counted unit: the Theorem whose source label is \texttt{thm:converse} in the article (source0.tex lines 699--720, source label \texttt{thm:converse}), delivered together with the article's own complete original proof of it (source0.tex lines 810--863, one \texttt{proof} environment of 54 lines). No mathematics is written, repaired or re-derived: statement and proof are the article's own text line for line. The proof is detached from the statement in the source: the article states the result at lines 699--720 and prints its proof at lines 810--863, and the proof environment's own header names the statement, so the association is the article's own. Required context retained uncounted: thm:gen\_gap\_2 (lines 724--739). Every definition, display and label of those lines is retained unchanged. Omitted from this package and not redistributed: 1--698, 721--723, 740--809, 864--994. Nothing inside a retained excerpt is omitted or altered: every retained body is delivered exactly as the article's lines are written, so no omission receipt of any kind accompanies this package. Citations: the retained excerpts carry no citation command at all, so no bibliography entry of the article is transcribed and no reference is redistributed. Reference adaptation: none was needed and none was made; every reference inside the delivered unit resolves inside it, so no unresolved reference or citation is produced. Printed slips kept literal: no printed word, symbol, hypothesis or number of the counted statement or of its proof was altered. Layout and class: the article's own class is \texttt{article} with packages including geometry, amssymb, amsmath, amsthm, pgfplots, pgfplotstable, arydshln, verbatim, titlesec, subfigure, algorithm2e, bm, pdfpages, enumerate; the wrapper is the lane's pinned \texttt{amsart} editorial wrapper at 10pt on A4, which restates the theorem-like environments the retained text needs and adds the article's own macro definitions that the retained text needs (none), with extra wrapper packages (bm, bbm, dsfont, xspace, cancel, mathtools, cleveref). The delivered numbering is the unit's own. Assets: no retained span contains a figure environment, a graphics-inclusion command, a PDF-page inclusion, a table or a TikZ picture; no image file is copied into this package, the package declares no asset dependency and no image file is redistributed with it. Licence: version arXiv:2410.11136v2 of this article is distributed under the Creative Commons Attribution 4.0 International licence, as recorded in the arXiv licence metadata for this exact version and shown on the abstract page https://arxiv.org/abs/2410.11136v2. A full rights notice is delivered at licenses/arxiv-2410.11136-CC-BY-4.0.txt. Characters: every retained excerpt is pure ASCII, so the delivered engine, XeLaTeX with its default Latin Modern text font, reports no missing character.
\begin{theorem}
\label{thm:gen_gap_2}
        Let $\mathcal{H}$ be a Hilbert space and suppose $\mathsf{P}_1,\ldots, \mathsf{P}_n$ are orthogonal projections with respect to the inner product of $\mathcal{H}$. Suppose that
        \begin{equation}
    \left\|\mathsf{P}_n\ldots \mathsf{P}_1\right\|_{\mathsf{op}}\leq 1-\delta.
\end{equation} 
Let $\mathsf{P}_{i_L}\ldots \mathsf{P}_{i_1}$ be any product of these orthogonal projection operators satisfying the following condition. For each $j\in [n]$, suppose there exists a minimal index $k_j\leq L$ satisfying $i_{k_j}=j$. We further assume these indices satisfy:
\begin{equation}
\label{eq:order_2}
        1=k_1\leq k_2\leq\ldots\leq k_n=L.
\end{equation}
Then, it holds that
\begin{equation*}
    \|\mathsf{P}_{i_L}\ldots \mathsf{P}_{i_1}\|_{\mathsf{op}}\leq 1-\frac{\delta^2}{8(L-n+1)}.
\end{equation*}
\end{theorem}

\begin{theorem}
\label{thm:converse}
    Let $\mathcal{H}$ be a Hilbert space and suppose $\mathsf{P}_1,\ldots,\mathsf{P}_n$ are orthogonal projections with respect to the inner product of $\mathcal{H}$. There is an absolute constant $c_{\ref{thm:converse}}>0$ such that the following holds for any $\delta>0$:
    \begin{enumerate}
        \item \label{item:all} Suppose that for every permutation $\sigma:[n]\to [n]$, it holds that
        \begin{equation*}
            \left\|\mathsf{P}_{\sigma(n)}\ldots \mathsf{P}_{\sigma(1)}\right\|_{\mathsf{op}}\leq 1-\delta.
        \end{equation*}
        Then 
        \begin{equation*}
            \left\|\frac{1}{n}\sum_{i=1}^n \mathsf{P}_i\right\|\leq 1-c_{\ref{thm:converse}}\left(\frac{\delta}{n\log(n/\delta)}\right)^2.
        \end{equation*}
        \item \label{item:one} Suppose that there exists a  permutation $\sigma:[n]\to [n]$ such that
        \begin{equation*}
            \left\|\mathsf{P}_{\sigma(n)}\ldots \mathsf{P}_{\sigma(1)}\right\|_{\mathsf{op}}\leq 1-\delta.
        \end{equation*}
        Then 
        \begin{equation*}
            \left\|\frac{1}{n}\sum_{i=1}^n \mathsf{P}_i\right\|\leq 1-c_{\ref{thm:converse}}\left(\frac{\delta}{n^2\log(n/\delta)}\right)^2.
        \end{equation*}
    \end{enumerate}
\end{theorem}

\begin{proof}[Proof of \Cref{thm:converse}]
    The idea is to take powers and argue that most of the terms contain one of the permutations as an embedded subsequence that is assumed to have the requisite gap. The spectral gap from the previous line will decay linearly in the power, but the probability of failing to contain the subsequence will decay like a large polynomial, so we can set parameters appropriately. We may assume $n\geq 2$ throughout.

    More formally, since by self-adjointness, we have
    \begin{equation*}
        \left\|\frac{1}{n}\sum_{i=1}^n \mathsf{P}_i\right\|^L =\left\|\left(\frac{1}{n}\sum_{i=1}^n \mathsf{P}_i\right)^L\right\| ,
    \end{equation*}
    it will suffice to upper bound the latter. For any $L\geq n$, we first write
    \begin{equation*}
        \left(\frac{1}{n}\sum_{i=1}^n \mathsf{P}_i \right)^L = \frac{1}{n^L}\sum_{i_1,\ldots,i_L=1}^n \mathsf{P}_{i_L}\ldots \mathsf{P}_{i_1}.
    \end{equation*}

    We now choose parameters for each case:
    \begin{itemize}
        \item First suppose the condition of \Cref{item:all} so that all scans have an operator norm gap. For any $L\geq n$, define 
        \begin{equation*}
            \mathcal{A} =\left\{(i_L,\ldots,i_1)\in [n]^L: [n] = \{i_L,\ldots,i_1\}\right\}.
        \end{equation*}
        In words, this is the set of all sequences of indices containing all indices in $[n]$. By standard coupon collector bounds, if we choose $L=5n \log(200n/\delta)$, then
        \begin{equation*}
            \Pr((i_L,\ldots,i_1)\not\in \mathcal{A})\leq \frac{\delta^{3}}{200n^{3}}.
        \end{equation*}
        On the other hand, for any such tuple $(i_L,\ldots,i_1)\in \mathcal{A}$, we claim
        \Cref{thm:gen_gap_2} implies that
        \begin{equation*}
            \|\mathsf{P}_{i_L}\ldots \mathsf{P}_{i_1}\|_{\mathsf{op}}\leq 1-\frac{\delta^2}{8L}=1-\frac{\delta^2}{40n\log(200n/\delta)}.
        \end{equation*}
        Indeed, this tuple contains all elements in some order by definition of $\mathcal{A}$, and we may assume $i_L=\sigma(n)$ is the last element found and $i_1=\sigma(1)$ is the first element of this permutation.  The reason is that the operator norm can only decrease by removing a prefix (on the left) of operators since projections cannot increase the norm and by removing any suffix of projectors (on the right) since these projectors map the unit ball of $\mathcal{H}$ to itself. But then by symmetry, we can directly apply \Cref{thm:gen_gap_2} to conclude the bound. It follows that
        \begin{align*}
            \left\|\left(\frac{1}{n}\sum_{i=1}^n \mathsf{P}_i\right)^L\right\|&\leq \frac{1}{n^L}\sum_{(i_L,\ldots,i_1)\in \mathcal{A}} \|\mathsf{P}_{i_L}\ldots \mathsf{P}_{i_1}\|_{\mathsf{op}}+\frac{1}{n^L}\sum_{(i_L,\ldots,i_1)\in \mathcal{A}^c} \|\mathsf{P}_{i_L}\ldots \mathsf{P}_{i_1}\|_{\mathsf{op}}\\
            &\leq 1-\frac{\delta^2}{40n\log(200n/\delta)}+\Pr((i_L,\ldots,i_1)\not\in \mathcal{A})\\
            &\leq 1-\frac{\delta^2}{40n\log(200n/\delta)}+\frac{\delta^{3}}{200n^3}\\
            &\leq 1-\frac{\delta^2}{100n\log(200n/\delta)}.
        \end{align*}
        Here, we simply apply the triangle inequality and crudely bound the operator norms by $1$ for those elements not in $\mathcal{A}$. Therefore,
        \begin{equation*}
            \left\|\frac{1}{n}\sum_{i=1}^n \mathsf{P}_i\right\|\leq \left(1-\frac{\delta^2}{100n\log(200n/\delta)}\right)^{1/L}\leq 1-\frac{c\delta^2}{n^2\log^2(n/\delta)}
        \end{equation*}
        for some absolute constant $c>0$.

        \item Next, we assume that there exists a permutation $\sigma:[n]\to [n]$ such that the associated product of the projections has operator norm at most $1-\delta$. We use the same approach as last time, now just with $L = Cn^2\log(100Cn/\delta)$ for a suitable constant $C>0$ to be determined later. Let $\mathcal{A}$ denote the set of tuples $(i_L,\ldots,i_1)$ containing $\sigma$ as an ordered subsequence. Note that the probability that $(i_L,\ldots,i_1)\not\in \mathcal{A}$ is precisely the probability that a $\mathsf{Bin}(L,1/n)$ random variable is less than $n$ by the natural coupling that interprets each success as sampling the next element of the permutation $\sigma$. For this choice of $L$, it is easy to see from Hoeffding's inequality that if $C$ is a large enough constant, then we have
        \begin{equation*}
            \Pr((i_L,\ldots,i_1)\not\in \mathcal{A})\leq \delta^{100}/100Cn^{100}.
        \end{equation*}
        The same argument as the previous case asserts that for any $(i_L,\ldots,i_1)\in \mathcal{A}$, we have an operator norm bound of $1-\delta^2/8Cn^2\log(100Cn/\delta)$. Therefore, the same argument and calculation splitting the sum across $\mathcal{A}$ yields
        \begin{align*}
            \left\|\left(\frac{1}{n}\sum_{i=1}^n \mathsf{P}_i\right)^L\right\|&\leq \frac{1}{n^L}\sum_{(i_L,\ldots,i_1)\in \mathcal{A}} \|\mathsf{P}_{i_L}\ldots \mathsf{P}_{i_1}\|_{\mathsf{op}}+\frac{1}{n^L}\sum_{(i_L,\ldots,i_1)\in \mathcal{A}^c} \|\mathsf{P}_{i_L}\ldots \mathsf{P}_{i_1}\|_{\mathsf{op}}\\
            &\leq 1-\frac{\delta^2}{8Cn^2\log(100Cn/\delta)}+\Pr((i_L,\ldots,i_1)\not\in \mathcal{A})\\
            &\leq 1-\frac{\delta^2}{8Cn^2\log(100Cn/\delta)}+\frac{\delta^{100}}{100Cn^{100}}\\
            &\leq 1-\frac{c\delta^2}{n^2\log(n/\delta)}.
        \end{align*}
        for a suitable constant $c>0$. Taking $1/L$ norms again gives the stated bound.\qedhere
    \end{itemize}
\end{proof}

\end{document}
